\section{Binární halda}\label{sec:halda}
Uveďme na úvod motivační příklad. Mějme seznam čísel,~z~něhož chceme co nejrychleji vybrat minimální nebo maximální hodnotu. Pro maximum bychom mohli sestavit následující jednoduchý algoritmus.
\begin{algorithm}
    \caption{Vyhledání maximální hodnoty v~seznamu (\textsc{Max})}
    \label{alg:max}
    \KwIn{Seznam čísel $x_1,x_2,\dots,x_n$}
    $m\gets x_1$\\
    \For{$i=2,3,\dots,n$}{
        \If{$x_i>m$}{
            $m\gets x_i$
        }
    }
    \Return{m}\\
    \KwOut{Maximální hodnota seznamu \textit{m}}
\end{algorithm}

Časová složitost bude zjevně $\bigO(n)$,~neboť algoritmus prochází všech $n$ prvků. Pokud ale se seznamem pracujeme opakovaně a~průběžně jej udržujeme vhodně uspořádaný,~můžeme minimum či maximum získat rychleji. K~tomu můžeme použít tzv. \emph{haldu}.
\begin{definition}[Minimová binární halda]\label{def:binarni_halda}
    Minimová binární halda je datová struktura tvaru binárního stromu s~množinou vrcholů $V$ a~klíčovou funkcí $\mapping{\key}{V}{\Z}$,~kde je v~každém vrcholu uložena \emph{právě jedna} hodnota $\key(v)$ (tzv. \emph{klíč}) a~navíc platí:
    \begin{enumerate}[label=(\roman*)]
        \item\label{binhalda_podminka_1} každá hladina je \emph{plně obsazena},~kromě poslední,~přičemž vrcholy poslední hladiny jsou zaplněny \emph{zleva},
        \item\label{binhalda_podminka_2} a~je-li $v$ libovolný vrchol a~$s$ jeho syn,~pak $\key(v)\leqslant\key(s)$.
    \end{enumerate}
\end{definition}

\begin{figure}[ht]
    \centering
    \begin{subfigure}{6cm}
        \centering
        \begin{tikzpicture}[x=1cm,y=1cm]
\node[vertex,minimum size=6mm] (1) at (0,0) {$2$};
\node[vertex,minimum size=6mm] (2) at (-1.5,-1.3) {$5$};
\node[vertex,minimum size=6mm] (3) at ( 1.5,-1.3) {$7$};
\node[vertex,minimum size=6mm] (4) at (-2.3,-2.6) {$9$};
\node[vertex,minimum size=6mm] (5) at (-0.7,-2.6) {$6$};
\node[vertex,minimum size=6mm] (6) at ( 0.7,-2.6) {$8$};
\node[vertex,minimum size=6mm] (7) at ( 2.3,-2.6) {$12$};

\draw[edge] (1)--(2);
\draw[edge] (1)--(3);
\draw[edge] (2)--(4);
\draw[edge] (2)--(5);
\draw[edge] (3)--(6);
\draw[edge] (3)--(7);
\end{tikzpicture}
        \caption{Korektní halda.}
        \label{subfig:korektni_halda}
    \end{subfigure}
    \quad
    \begin{subfigure}{5cm}
        \centering
        \begin{tikzpicture}[x=1cm,y=1cm]
\node[vertex,minimum size=6mm] (1) at (0,0) {$1$};
\node[vertex,minimum size=6mm] (2) at (-1.5,-1.3) {$4$};
\node[vertex,minimum size=6mm] (3) at ( 1.5,-1.3) {$6$};
\node[vertex,minimum size=6mm,color=red] (5) at ( 0.7,-2.6) {$9$};

\draw[edge] (1)--(2);
\draw[edge] (1)--(3);
\draw[edge] (3)--(5);
\end{tikzpicture}
        \caption{Mezera před obsazeným místem v~poslední hladině.}
        \label{subfig:nekorektni_halda_1}
    \end{subfigure}
    \quad
    \begin{subfigure}{5cm}
        \centering
        \begin{tikzpicture}[x=1cm,y=1cm]
    \node[vertex,minimum size=6mm,draw=red] (r) at (0,0) {$2$};
    \node[vertex,minimum size=6mm,draw=red] (l) at (-1.4,-1.3) {$1$};
    \node[vertex,minimum size=6mm] (p) at (1.4,-1.3) {$5$};
    \node[vertex,minimum size=6mm] (a) at (-2,-2.6) {$7$};
    \node[vertex,minimum size=6mm] (b) at (-0.7,-2.6) {$8$};
    \node[vertex,minimum size=6mm] (c) at (0.7,-2.6) {$6$};
    \node[vertex,minimum size=6mm] (d) at (2,-2.6) {$9$};
    \foreach \x/\y in {r/l,r/p,l/a,l/b,p/c,p/d}\draw[edge](\x)--(\y);
\end{tikzpicture}

        \caption{Klíč levého syna je menší než klíč kořene.}
        \label{subfig:nekorektni_halda_2}
    \end{subfigure}
    \caption{Příklady korektních a~nekorektních hald.}
    \label{fig:halda_modely}
\end{figure}

Zde se na chvíli pozastavme. Podmínka \ref{binhalda_podminka_2} v~definici binární haldy \ref{def:binarni_halda} má velmi příjemný důsledek. Pokud se budeme pohybovat od kořene směrem dolů,~hodnoty klíčů se nemohou zmenšovat. Minimum se proto vždy nachází \emph{v~kořeni stromu},~takže jeho zjištění zvládneme v~\emph{konstantním čase} $\bigO(1)$.
\notebox{Pokud bychom chtěli naopak \emph{maximovou binární haldu},~bude definice \ref{def:binarni_halda} vypadat obdobně,~akorát v~podmínce \ref{binhalda_podminka_2} bude obrácená nerovnost,~tj. $\key(v)\geqslant\key(s)$ (klíč v~rodiči má vždy stejnou nebo vyšší hodnotu než klíče v~jeho synech). Maximum se bude opět nacházet v~kořeni haldy.}
Je však více operací,~které bychom rádi s~haldou prováděli. Vypišme si všechny,~které nás budou zajímat.
\begin{itemize}
    \item \textsc{Insert}($x$) -- \emph{vložení} nového vrcholu s~klíčem $x$ do haldy.
    \item \textsc{Min}() -- \emph{nalezení} vrcholu s~nejmenším klíčem (už jsme zmínili).
    \item \textsc{ExtractMin}() -- \emph{odebrání} vrcholu s~nejmenším klíčem.
    \item \textsc{Increase}($v$) -- \emph{zvýšení} hodnoty klíče ve zvoleném vrcholu $v$.
    \item \textsc{Decrease}($v$) -- \emph{snížení} hodnoty klíče ve zvoleném vrcholu $v$.
\end{itemize}

\subsection{Halda uložená v~poli}

Úplný tvar binární haldy dovoluje vynechat ukazatele mezi vrcholy a~uložit klíče přímo do pole. Vrcholy očíslujeme po hladinách zleva doprava,~počínaje kořenem s~indexem $1$. Je-li vrchol na pozici $i$,~pak jeho levý syn má index $2i$,~pravý syn $2i+1$ a~rodič (pro $i>1$) index $\lfloor i/2\rfloor$.

Protože jsou všechny hladiny kromě poslední plné a~poslední se zaplňuje zleva,~nevznikají v~poli žádná prázdná místa. První volná pozice je vždy konec pole,~takže vložení nového klíče i~odebrání posledního vrcholu jsou jednoduché. Samotné obnovení haldového uspořádání pak probíhá prohazováním prvků na uvedených indexech.

V~programovacích jazycích s~indexací pole od nuly se vzorce posunou: děti prvku $i$ leží na pozicích $2i+1$ a~$2i+2$,~rodič na $\lfloor(i-1)/2\rfloor$. Jde o~stejnou datovou strukturu; před implementací je pouze nutné zvolit jednu konvenci a~důsledně ji dodržet.

Podívejme se postupně na jednotlivé operace a~jejich realizaci. Slovní popis vždy doplníme pseudokódem. V~něm budeme haldu zapisovat jako pole $H[1],\dots,H[n]$ indexované od jedničky; proměnná $n$ označuje jeho aktuální délku.

\subsection{Vkládání nového vrcholu}

Při vkládání nového vrcholu (\textsc{Insert}) nejdříve potřebujeme rozhodnout,~kam jej umístíme. S~tím nám pomůže podmínka \ref{binhalda_podminka_1}: nový vrchol vložíme na první volné místo v~pořadí zleva doprava. Pokud je dosavadní poslední hladina plná,~založíme novou hladinu a~vrchol umístíme zcela vlevo. Tím zachováme tvar haldy,~ale nemusí být splněna podmínka \ref{binhalda_podminka_2}. Klíč nového vrcholu může být menší než klíč jeho rodiče. Haldu opravíme postupným prohazováním klíče s~rodičem,~dokud nebude podmínka splněna (viz obrázek \ref{fig:vkladani_vrcholu_halda}).
\begin{figure}[ht]
    \centering
    \begin{subfigure}{7cm}
        \input{01-grafalgo/images/ch01_vkladani_1.tex}
    \end{subfigure}
    \begin{subfigure}{7cm}
        \begin{tikzpicture}[scale=0.85]
    \node[vertex] (a) at (0,0) {$1$};
    \node[vertex] (b) at (-1.5,-1.3) {$4$};
    \node[vertex,draw=red] (c) at (1.5,-1.3) {$2$};
    \node[vertex] (d) at (-2.2,-2.6) {$7$};
    \node[vertex] (e) at (-0.8,-2.6) {$8$};
    \node[vertex] (f) at (0.8,-2.6) {$6$};
    \foreach \x/\y in {a/b,a/c,b/d,b/e,c/f}\draw[edge](\x)--(\y);
\end{tikzpicture}

    \end{subfigure}
    \caption{Vkládání nového vrcholu do haldy.}
    \label{fig:vkladani_vrcholu_halda}
\end{figure}
Postup můžeme zapsat zkráceně ve třech bodech takto:
\begin{itemize}
    \item Vložíme nový vrchol s~klíčem $x$ na první volnou pozici zleva na poslední hladině.
    \item pokud je klíč rodiče větší než $x$,~prohodíme vrcholy (resp. jejich klíče)\footnote{Vrchol v~podstatě takto "probublá" do správné pozice ve stromě (příp. až do pozice kořene).}.
    \item Opakujeme druhý krok.
\end{itemize}

Samotné probublávání nového klíče směrem ke kořeni oddělíme do pomocné operace \textsc{BubbleUp}. Díky tomu ji později použijeme beze změny také při snížení existujícího klíče.
\begin{algorithm}
    \caption{\textsc{bubble-up}}
    \label{alg:heap-bubble-up}
    \KwIn{Minimová halda $H[1],\dots,H[n]$ a~index $i\in\set{1,\dots,n}$}
    \While{$i>1$ \textup{a} $H[\lfloor i/2\rfloor]>H[i]$}{
        $r\gets\lfloor i/2\rfloor$\tcp*{index rodiče}
        Prohoď $H[i]$ a~$H[r]$\;
        $i\gets r$\;
    }
\end{algorithm}

\subsection{Odstranění minima}

Získání minimálního klíče v~haldě je triviální: přečteme klíč v~kořeni a~jsme hotovi. Pokud však chceme minimum operací \textsc{ExtractMin} také odstranit,~je to o~něco složitější. Kořen nemůžeme jen smazat,~jeho pozici musí zastoupit jiný vrchol. Abychom neporušili tvar haldy,~přesuneme do kořene \emph{poslední vrchol v~pořadí zleva doprava},~tedy nejpravější obsazený vrchol poslední hladiny,~a~jeho původní pozici odstraníme.

Nyní však (stejně jako u~vkládání vrcholu) nastává problém s~druhou podmínkou,~neboť touto operací jsme ji mohli porušit. Provedeme obdobný proces,~jako při vkládání vrcholu do haldy,~ale vrchol nyní bude "probublávat" směrem dolů. Podíváme se na syny daného vrcholu a~pokud klíč některého z~nich je menší než klíč nového kořene,~pak jej s~ním prohodíme (v~případě,~že klíče obou synů jsou menší,~než klíč kořene,~prohodíme s~menším z~nich).
\begin{figure}[ht]
    \centering
    \begin{subfigure}{5cm}
        \begin{tikzpicture}[scale=0.7]
    \node[vertex,draw=red] (a) at (0,0) {$1$};
    \node[vertex] (b) at (-1.5,-1.3) {$3$};
    \node[vertex] (c) at (1.5,-1.3) {$4$};
    \node[vertex] (d) at (-2.2,-2.6) {$7$};
    \node[vertex] (e) at (-0.8,-2.6) {$8$};
    \node[vertex] (f) at (0.8,-2.6) {$6$};
    \node[vertex,draw=red] (g) at (2.2,-2.6) {$9$};
    \foreach \x/\y in {a/b,a/c,b/d,b/e,c/f,c/g}\draw[edge](\x)--(\y);
    \draw[red] ($(a)+(-.25,.25)$)--($(a)+(.25,-.25)$);
\end{tikzpicture}

    \end{subfigure}
    \begin{subfigure}{5cm}
        \input{01-grafalgo/images/ch01_odebirani_2.tex}
    \end{subfigure}
    \begin{subfigure}{5cm}
        \input{01-grafalgo/images/ch01_odebirani_3.tex}
    \end{subfigure}
    \begin{subfigure}{5cm}
        \begin{tikzpicture}[scale=0.7]
    \node[vertex] (a) at (0,0) {$3$};
    \node[vertex,draw=red] (b) at (-1.5,-1.3) {$9$};
    \node[vertex] (c) at (1.5,-1.3) {$4$};
    \node[vertex,draw=red] (d) at (-2.2,-2.6) {$7$};
    \node[vertex] (e) at (-0.8,-2.6) {$8$};
    \node[vertex] (f) at (0.8,-2.6) {$6$};
    \foreach \x/\y in {a/b,a/c,b/d,b/e,c/f}\draw[edge](\x)--(\y);
    \draw[red,diredge,bend right=45] (b) to (d);
\end{tikzpicture}

    \end{subfigure}
    \begin{subfigure}{5cm}
        \begin{tikzpicture}[scale=0.7]
    \node[vertex] (a) at (0,0) {$3$};
    \node[vertex] (b) at (-1.5,-1.3) {$7$};
    \node[vertex] (c) at (1.5,-1.3) {$4$};
    \node[vertex] (d) at (-2.2,-2.6) {$9$};
    \node[vertex] (e) at (-0.8,-2.6) {$8$};
    \node[vertex] (f) at (0.8,-2.6) {$6$};
    \foreach \x/\y in {a/b,a/c,b/d,b/e,c/f}\draw[edge](\x)--(\y);
\end{tikzpicture}

    \end{subfigure}
    \caption{Odebírání vrcholu (kořene) s~minimálním klíčem.}
    \label{fig:odebirani_vrcholu_halda}
\end{figure}
V~bodech můžeme zapsat celou operaci následovně:
\begin{itemize}
    \item Smažeme kořen a~nahradíme jej nejpravějším obsazeným vrcholem na poslední hladině.
    \item Tento vrchol (resp. jeho klíč) porovnáme s~jeho syny a~případně prohodíme s~tím,~jehož klíč je menší.
    \item Opakujeme druhý krok.
\end{itemize}
Operace,~při nichž vrchol během vkládání \uv{probublává} nahoru nebo během odebírání minima dolů,~se někdy nazývají \textsc{BubbleUp} a~\textsc{BubbleDown}. Ještě se nám budou hodit u~operací \textsc{Increase} a~\textsc{Decrease}.

Při probublávání dolů musíme v~každém kroku vybrat menšího ze synů. Prohození s~větším synem by nemuselo obnovit haldovou podmínku vůči tomu menšímu. Pokud pravý syn neexistuje, porovnáváme pouze s~levým.
\begin{algorithm}
    \caption{\textsc{bubble-down}}
    \label{alg:heap-bubble-down}
    \KwIn{Minimová halda $H[1],\dots,H[n]$ a~index $i\in\set{1,\dots,n}$}
    \While{$2i\leqslant n$}{
        $s\gets 2i$\tcp*{zatím levý syn}
        \If{$2i+1\leqslant n$ \textup{a} $H[2i+1]<H[s]$}{
            $s\gets 2i+1$\tcp*{pravý syn je menší}
        }
        \If{$H[i]\leqslant H[s]$}{
            \Return\;
        }
        Prohoď $H[i]$ a~$H[s]$\;
        $i\gets s$\;
    }
\end{algorithm}

Pomocí této operace už lze \textsc{ExtractMin} zapsat velmi stručně. Zvlášť ošetřujeme jednoprvkovou haldu: po odebrání posledního prvku již není co probublávat.
\begin{algorithm}
    \caption{\textsc{extract-min}}
    \label{alg:heap-extract-min}
    \KwIn{Neprázdná minimová halda $H[1],\dots,H[n]$}
    $m\gets H[1]$\;
    $H[1]\gets H[n]$\;
    Odstraň poslední prvek pole $H$\;
    $n\gets n-1$\;
    \If{$n>0$}{
        Zavolej \textsc{bubble-down}$(H,1)$\;
    }
    \Return{$m$}\;
    \KwOut{Odebraný minimální klíč $m$}
\end{algorithm}

\subsection{Zvýšení a~snížení hodnoty klíče ve vrcholu}

Poslední dvojicí operací,~která se nám bude hodit,~je zvýšení (\textsc{Increase}) a~snížení (\textsc{Decrease}) klíče ve vybraném vrcholu. Předpokládáme přitom,~že umíme daný vrchol v~haldě přímo najít. Pokud v~haldě uchováváme složitější objekty,~můžeme si například vést pomocný slovník,~který každému objektu přiřadí jeho současnou pozici v~haldě. Při každém prohození pak aktualizujeme i~slovník. Ten spotřebuje navíc $\bigO(n)$ paměti.

Realizace samotných operací je již jednoduchá. Pokud provedeme \textsc{Increase} klíče v~určitém vrcholu,~může se halda pokazit pouze směrem dolů. Provedeme tedy \textsc{BubbleDown},~stejně jako při odebírání kořene.

\begin{figure}[ht]
    \centering
    \begin{tikzpicture}[scale=0.82]
    \begin{scope}[xshift=-4.4cm]
        \node[above] at (0,0.75) {po zvýšení $4\to 9$};
        \node[vertex] (a1) at (0,0) {$1$};
        \node[vertex,draw=red] (b1) at (-1.5,-1.3) {$9$};
        \node[vertex] (c1) at (1.5,-1.3) {$3$};
        \node[vertex,draw=red] (d1) at (-2.2,-2.6) {$7$};
        \node[vertex] (e1) at (-0.8,-2.6) {$8$};
        \node[vertex] (f1) at (0.8,-2.6) {$6$};
        \node[vertex] (g1) at (2.2,-2.6) {$5$};
        \foreach \x/\y in {a1/b1,a1/c1,b1/d1,b1/e1,c1/f1,c1/g1}
            \draw[edge] (\x)--(\y);
        \draw[red,diredge,bend right=38] (b1) to (d1);
    \end{scope}

    \draw[diredge] (-0.75,-1.3) -- (0.75,-1.3)
        node[midway,above] {\textsc{bubble-down}};

    \begin{scope}[xshift=4.4cm]
        \node[above] at (0,0.75) {obnovená halda};
        \node[vertex] (a2) at (0,0) {$1$};
        \node[vertex,draw=red] (b2) at (-1.5,-1.3) {$7$};
        \node[vertex] (c2) at (1.5,-1.3) {$3$};
        \node[vertex,draw=red] (d2) at (-2.2,-2.6) {$9$};
        \node[vertex] (e2) at (-0.8,-2.6) {$8$};
        \node[vertex] (f2) at (0.8,-2.6) {$6$};
        \node[vertex] (g2) at (2.2,-2.6) {$5$};
        \foreach \x/\y in {a2/b2,a2/c2,b2/d2,b2/e2,c2/f2,c2/g2}
            \draw[edge] (\x)--(\y);
    \end{scope}
\end{tikzpicture}

    \caption{Operace \textsc{Increase}: zvýšený klíč se pomocí \textsc{BubbleDown} přesune dolů.}
    \label{fig:halda-increase}
\end{figure}

Operaci \textsc{Decrease} provedeme zcela stejně,~akorát nyní se halda bude kazit "směrem nahoru",~takže provedeme \textsc{BubbleUp},~jako při vkládání nového vrcholu (\textsc{Insert}).

\begin{figure}[ht]
    \centering
    \begin{tikzpicture}[scale=0.82]
    \begin{scope}[xshift=-4.4cm]
        \node[above] at (0,0.75) {po snížení $6\to 2$};
        \node[vertex] (a1) at (0,0) {$1$};
        \node[vertex] (b1) at (-1.5,-1.3) {$4$};
        \node[vertex,draw=red] (c1) at (1.5,-1.3) {$3$};
        \node[vertex] (d1) at (-2.2,-2.6) {$7$};
        \node[vertex] (e1) at (-0.8,-2.6) {$8$};
        \node[vertex,draw=red] (f1) at (0.8,-2.6) {$2$};
        \node[vertex] (g1) at (2.2,-2.6) {$5$};
        \foreach \x/\y in {a1/b1,a1/c1,b1/d1,b1/e1,c1/f1,c1/g1}
            \draw[edge] (\x)--(\y);
        \draw[red,diredge,bend left=38] (f1) to (c1);
    \end{scope}

    \draw[diredge] (-0.75,-1.3) -- (0.75,-1.3)
        node[midway,above] {\textsc{bubble-up}};

    \begin{scope}[xshift=4.4cm]
        \node[above] at (0,0.75) {obnovená halda};
        \node[vertex] (a2) at (0,0) {$1$};
        \node[vertex] (b2) at (-1.5,-1.3) {$4$};
        \node[vertex,draw=red] (c2) at (1.5,-1.3) {$2$};
        \node[vertex] (d2) at (-2.2,-2.6) {$7$};
        \node[vertex] (e2) at (-0.8,-2.6) {$8$};
        \node[vertex,draw=red] (f2) at (0.8,-2.6) {$3$};
        \node[vertex] (g2) at (2.2,-2.6) {$5$};
        \foreach \x/\y in {a2/b2,a2/c2,b2/d2,b2/e2,c2/f2,c2/g2}
            \draw[edge] (\x)--(\y);
    \end{scope}
\end{tikzpicture}

    \caption{Operace \textsc{Decrease}: snížený klíč se pomocí \textsc{BubbleUp} přesune nahoru.}
    \label{fig:halda-decrease}
\end{figure}

Pseudokód obou operací zároveň kontroluje, že nový klíč skutečně odpovídá názvu operace. Kdybychom například v~operaci \textsc{Increase} předali menší hodnotu, bylo by potřeba opravovat haldu opačným směrem.
\begin{algorithm}
    \caption{\textsc{increase}}
    \label{alg:heap-increase}
    \KwIn{Minimová halda $H[1],\dots,H[n]$, index $i$ a~nový klíč $x$}
    \If{$x<H[i]$}{
        Ohlas chybu\;
    }
    $H[i]\gets x$\;
    Zavolej \textsc{bubble-down}$(H,i)$\;
\end{algorithm}

\begin{algorithm}
    \caption{\textsc{decrease}}
    \label{alg:heap-decrease}
    \KwIn{Minimová halda $H[1],\dots,H[n]$, index $i$ a~nový klíč $x$}
    \If{$x>H[i]$}{
        Ohlas chybu\;
    }
    $H[i]\gets x$\;
    Zavolej \textsc{bubble-up}$(H,i)$\;
\end{algorithm}
\notebox{V~případě \emph{maximové binární haldy} by operace vypadaly zcela stejně,~akorát při operaci \textsc{Increase} budeme opravovat haldu směrem nahoru a~u~\textsc{Decrease} směrem dolů.}

\subsection{Implementace v~Pythonu}

Program~\ref{lst:min-heap-python} ukládá minimovou haldu v~seznamu indexovaném od nuly. Proto používá vztahy $2i+1$ a~$2i+2$ pro syny a~$\lfloor(i-1)/2\rfloor$ pro rodiče. Metody \texttt{bubble\_up} a~\texttt{bubble\_down} přesně odpovídají pomocným operacím z~pseudokódu; liší se pouze posunutými indexy.

Konstruktor vytváří haldu opakovaným vkládáním, což je didakticky přímočaré, i~když pro velké pole existuje rychlejší hromadná konstrukce. Metody \texttt{increase} a~\texttt{decrease} očekávají index prvku. V~praktickém programu bychom při ukládání složitějších objektů přidali pomocný slovník objekt--index a~aktualizovali jej při každém prohození.

\lstinputlisting[
    caption={Minimová binární halda a~její základní operace v~Pythonu.},
    label={lst:min-heap-python}
]{01-grafalgo/code/min_heap.py}

\subsection{Časové složitosti operací}

Pojďme si nyní rozebrat časové složitosti zmíněných operací \textsc{Insert},~\textsc{Min},~\textsc{ExtractMin},~\textsc{Increase} a~\textsc{Decrease} a~podívat se,~jak moc jsme si pomohli oproti práci s~polem (seznamem). Pokud se pozorně podíváme na operace popsané výše,~je zjevné,~že nejvíce bude naše operace zpomalovat ono "probublávání" vrcholu při opravování haldy (vyjma operace \textsc{Min},~kde jen přečteme klíč v~kořeni),~tj. \textsc{BubbleUp} a~\textsc{BubbleDown}. Na nich bude celková časová složitost záviset.
\tipbox{Všimněme si,~že po každém prohození vrcholů při \textsc{BubbleUp} nebo \textsc{BubbleDown} se vrchol posune o~jednu hladinu výše/níže.}
Víme tak,~že při opravování haldy se vrchol může posunout \emph{maximálně o~počet hladin},~který má daná halda. Časová složitost daných operací bude tak závislá na jejich počtu. Tím jsme se dostali o~trochu blíže řešení,~ale rádi bychom věděli,~jak se bude době trvání daných operací měnit v~závislosti na \emph{počtu vrcholů v~haldě},~nikoliv počtu hladin. Mezi těmito proměnnými však existuje hezká závislost.

Dejme tomu,~že naše halda má $n$ vrcholů v~celkově $h$ hladinách. Halda s~jednou plnou hladinou má jeden vrchol,~se dvěma plnými hladinami tři vrcholy,~se třemi sedm vrcholů atd. (viz tabulka \ref{tab:halda_vrcholy_hladiny}).
\begin{table}[ht]
    \centering
    \begin{tabular}{|l|c|c|c|c|c|c|}
        \hline
        Počet hladin $h$  & 1 & 2 & 3 & 4  & 5  & 6  \\ \hline
        Počet vrcholů $n$ & 1 & 3 & 7 & 15 & 31 & 63 \\ \hline
    \end{tabular}
    \caption{Vztah mezi počtem hladin $h$ a~počtem vrcholů $n$ v~plně obsazené haldě.}
    \label{tab:halda_vrcholy_hladiny}
\end{table}
\begin{figure}[ht]
    \centering
    \begin{tikzpicture}[x=1cm,y=1cm]

% souřadnice hladin
\def\yzero{0}
\def\yone{-1.4}
\def\ytwo{-2.8}
\def\ydots{-4.6}
\def\yk{-6.0}

% vrcholy
\node[vertex] (r) at (0,\yzero) {};

\node[vertex] (a1) at (-2,\yone) {};
\node[vertex] (a2) at ( 2,\yone) {};

\node[vertex] (b1) at (-3,\ytwo) {};
\node[vertex] (b2) at (-1,\ytwo) {};
\node[vertex] (b3) at ( 1,\ytwo) {};
\node[vertex] (b4) at ( 3,\ytwo) {};

\node[vertexplain] (d1) at (-3.6,\ydots) {$\vdots$};
\node[vertexplain] (d2) at (-1.2,\ydots) {$\vdots$};
\node[vertexplain] (d3) at ( 1.2,\ydots) {$\vdots$};
\node[vertexplain] (d4) at ( 3.6,\ydots) {$\vdots$};

\node[vertex] (k1) at (-4.2,\yk) {};
\node[vertex] (k2) at (-3.0,\yk) {};
\node[vertexplain] (kd1) at (-1.4,\yk) {$\cdots$};
\node[vertex] (k3) at (0,\yk) {};
\node[vertexplain] (kd2) at (1.4,\yk) {$\cdots$};
\node[vertex] (k4) at (3.0,\yk) {};
\node[vertex] (k5) at (4.2,\yk) {};

% hrany
\draw[edge] (r) -- (a1);
\draw[edge] (r) -- (a2);

\draw[edge] (a1) -- (b1);
\draw[edge] (a1) -- (b2);
\draw[edge] (a2) -- (b3);
\draw[edge] (a2) -- (b4);

% naznačení pokračování
\draw[edge, dashed] (b1) -- (-3.6,\ydots+0.4);
\draw[edge, dashed] (b2) -- (-1.2,\ydots+0.4);
\draw[edge, dashed] (b3) -- ( 1.2,\ydots+0.4);
\draw[edge, dashed] (b4) -- ( 3.6,\ydots+0.4);

\draw[edge, dashed] (-3.6,\ydots-0.25) -- (k1);
\draw[edge, dashed] (-3.6,\ydots-0.25) -- (k2);

\draw[edge, dashed] (-1.2,\ydots-0.25) -- ++(-0.3,-0.9);
\draw[edge, dashed] (-1.2,\ydots-0.25) -- ++( 0.3,-0.9);

\draw[edge, dashed] ( 1.2,\ydots-0.25) -- ++(-0.3,-0.9);
\draw[edge, dashed] ( 1.2,\ydots-0.25) -- ++( 0.3,-0.9);

\draw[edge, dashed] ( 3.6,\ydots-0.25) -- (k4);
\draw[edge, dashed] ( 3.6,\ydots-0.25) -- (k5);

% popisky hladin vlevo
\node[left] at (-5.2,\yzero) {$0$};
\node[left] at (-5.2,\yone)  {$1$};
\node[left] at (-5.2,\ytwo)  {$2$};
\node[left] at (-5.2,\ydots) {$\vdots$};
\node[left] at (-5.2,\yk)    {$k$};

% svislá závorka pro hladiny
\draw[decorate, decoration={brace, amplitude=5pt}]
    (-4.8,\yk) -- (-4.8,\yzero);

\end{tikzpicture}
    \caption{Plná binární halda s~obecně $h$ hladinami.}
    \label{fig:halda_hladiny}
\end{figure}
Je-li všech $h$ hladin plných,~platí $n=2^h-1$. V~obecné haldě jsou první $h-1$ hladiny plné a~poslední obsahuje alespoň jeden vrchol. Proto
\[2^{h-1}\leqslant n\leqslant 2^h-1.\]
Z~levé nerovnosti dostáváme $h\leqslant\log_2 n+1$,~tedy počet hladin je $\bigO(\log n)$. Tím máme vše připravené k~odvození složitosti operací.
\begin{theorem}[Složitost operací v~binární haldě]
    Operace \textsc{Insert},~\textsc{ExtractMin},~\textsc{Increase} a~\textsc{Decrease} mají časovou složitost $\bigO(\log{n})$. Operace \textsc{Min} má časovou složitost $\bigO(1)$.
\end{theorem}
\begin{proof}
    V~případě \textsc{Min} jen přečteme klíč v~kořeni,~což potrvá $\bigO(1)$. U~zbývajících operací opravujeme haldu pomocí \textsc{BubbleUp} nebo \textsc{BubbleDown}. Při každém prohození se zpracovávaný klíč posune o~jednu hladinu. Provede se tedy nejvýše $h-1$ prohození. Protože $h\leqslant\log_2 n+1$,~dostáváme
    \[\bigO(h)=\bigO(\log n).\]
    Operace \textsc{Insert},~\textsc{ExtractMin},~\textsc{Increase} a~\textsc{Decrease} proto mají logaritmickou časovou složitost.
\end{proof}
Zkusme si na závěr udělat menší porovnání oproti poli,~se kterým jsme začínali (viz tabulka \ref{tab:halda_pole_operace}).
\begin{table}[ht]
    \centering
    \begin{tabular}{l|ccccc}
                  & \textsc{Insert} & \textsc{Min} & \textsc{ExtractMin} & \textsc{Increase} & \textsc{Decrease} \\\hline
    Pole (seznam) & $\bigO(1)$                       & $\bigO(n)$                    & $\bigO(n)$                           & $\bigO(1)$                         & $\bigO(1)$                         \\
    Binární halda & $\bigO(\log{n})$                 & $\bigO(1)$                    & $\bigO(\log{n})$                     & $\bigO(\log{n})$                   & $\bigO(\log{n})$                  
    \end{tabular}
    \caption{Porovnání časových složitostí operací v~haldě a~v~poli.}
    \label{tab:halda_pole_operace}
\end{table}
Logaritmická časová složitost je určitě lepší,~než lineární,~je však vidět,~že jsme si zároveň pohoršili v~případě vkládání a~úpravy hodnot klíčů. Je tomu tak z~důvodu,~že pole má svou výhodu v~téměř nulové režii,~kdežto halda (kvůli své pevně definované struktuře) spotřebuje nějaký čas,~aby byla uvedena do korektního stavu.
